% BAB III - CARA PENGGUNAAN
\section{Cara Penggunaan}

\subsection{Menjalankan Aplikasi}

Untuk menjalankan aplikasi Gestune, ikuti langkah-langkah berikut:

\subsubsection{Jalankan Aplikasi Utama}
Jalankan file \texttt{main.py} dengan perintah:

\begin{lstlisting}[language=bash]
python main.py
\end{lstlisting}

Aplikasi juga dapat dijalankan dengan membuka run.bat yang telah disediakan (khusus Windows).

\subsubsection{Pengaturan Webcam}
\begin{itemize}
    \item Pastikan webcam Anda aktif dan terhubung dengan baik
    \item Posisikan diri Anda di depan webcam dengan pencahayaan yang cukup dengan jarak yang ideal (sekitar 60 - 100 cm)
    \item Pastikan kedua tangan dapat terlihat jelas dalam frame kamera
\end{itemize}

\subsection{Kontrol Aplikasi}
\subsubsection{Kontrol Panel GUI}
Setelah aplikasi dijalankan, pada bagian kanan layar akan muncul panel GUI yang berisi beberapa kontrol berikut:
\begin{itemize}
    \item \textbf{BPM Slider}: Slider untuk mengatur BPM (Beats Per Minute) dari pola ritme drum, tombol reset BPM ke default (Reset) serta tombol menambahkan dan mengurangi BPM (inkremental/dekremental 10).
    \item \textbf{Pattern Selector}: Tombol tombol untuk memilih pola ritme drum yang diinginkan
    \item \textbf{Visualisasi Instrumen}: Area visualisasi yang menampilkan instrumen drum yang sedang aktif
    \item \textbf{Exit}: Tombol untuk keluar dari aplikasi
\end{itemize}

\subsubsection{Tangan Kiri - Arpeggiator}
Tangan kiri digunakan untuk mengontrol Arpeggiator dengan dua cara:

\begin{enumerate}
    \item \textbf{Mengatur Pitch Nada}
    \begin{itemize}
        \item Naikkan tangan untuk mendapatkan nada yang lebih tinggi
        \item Turunkan tangan untuk mendapatkan nada yang lebih rendah
        \item Pitch nada akan berubah secara gradual mengikuti posisi vertikal tangan
    \end{itemize}
    
    \item \textbf{Mengatur Volume}
    \begin{itemize}
        \item Lakukan gesture pinch dengan mendekatkan ibu jari dan telunjuk
        \item Semakin kecil jarak antara ibu jari dan telunjuk, semakin besar volume
        \item Semakin besar jarak antara ibu jari dan telunjuk, semakin kecil volume
    \end{itemize}
\end{enumerate}

\subsubsection{Tangan Kanan - Drum Machine}
Tangan kanan digunakan untuk mengontrol Drum Machine dengan mengaktifkan instrumen drum berdasarkan jari yang diangkat:

\begin{table}[H]
\centering
\caption{Mapping Jari Drum Machine}
\label{tab:drum-mapping}
\begin{tabular}{lll}
\hline
\textbf{Jari} & \textbf{Instrumen} \\ \hline
Ibu Jari (Thumb) & Kick \\
Telunjuk (Index) & Snare \\
Tengah (Middle) & Hi-hat \\
Manis (Ring) & Tom \\
Kelingking (Pinky) & Crash \\
\hline
\end{tabular}
\end{table}

Setiap jari yang diangkat akan mengaktifkan track instrumen tersebut dalam loop sequencer yang sedang berjalan.

\paragraph{Penggantian Pola (Pattern Switching)}
Terdapat 7 pola ritme bawaan yang dapat dipilih pengguna:
\begin{itemize}
    \item Pattern 1: Modern Pop Groove
    \item Pattern 2: Dark Trap Groove
    \item Pattern 3: Minimal Trap Beat
    \item Pattern 4: Uptempo Breakbeat
    \item Pattern 5: Bounce Groove
    \item Pattern 6: Percussive Groove
    \item Pattern 7: Dance/Pop Rhtyhm
\end{itemize}

Untuk mengganti pola, pengguna dapat melakukan gesture \textbf{Kepalan Tangan (Fist)} pada tangan kanan. Sistem akan berpindah ke pola berikutnya secara siklis.

Pengguna juga dapat menambahkan pattern mereka sendiri dengan mengedit file patterns.py. 

\subsection{Tips Penggunaan}

Untuk mendapatkan pengalaman terbaik saat menggunakan aplikasi Gestune, perhatikan tips berikut:

\begin{enumerate}
    \item \textbf{Pencahayaan}: Pastikan ruangan memiliki pencahayaan yang cukup agar hand tracking dapat mendeteksi tangan dengan optimal
    
    \item \textbf{Jarak dengan Webcam}: Jaga jarak yang nyaman dengan webcam (sekitar 60-100 cm)
    
    \item \textbf{Posisi Tangan}: Pastikan kedua tangan terlihat jelas dalam frame kamera dan tidak tertutup oleh objek lain
    
    \item \textbf{Gerakan}: Lakukan gerakan dengan smooth dan tidak terlalu cepat agar sistem dapat mendeteksi dengan akurat
    
\end{enumerate}

\subsection{Troubleshooting}

Jika mengalami masalah saat menggunakan aplikasi, coba solusi berikut:

\begin{itemize}
    \item \textbf{Tangan tidak terdeteksi}: Perbaiki pencahayaan atau sesuaikan jarak dengan webcam
    \item \textbf{Audio tidak keluar}: Periksa pengaturan audio sistem Anda atau jalankan \texttt{test\_patterns.py} dan \textit{test\_audio\_engine.py} untuk verifikasi
    \item \textbf{Aplikasi lambat}: Tutup aplikasi lain yang berjalan untuk mengalokasikan lebih banyak resource
\end{itemize}
